\documentclass[11pt]{article}
\usepackage{amsmath,hyperref}
\title{Do not route a withdrawal through the remaining sources}
\date{}
\begin{document}
\maketitle

A withdrawn carrier can stop on the only shelf from which its remaining order can be completed. The \href{https://github.com/eternal-november/eos-research-programmes/pull/415}{notice experiment} on pinned MRWS found this for an earlier recovery rule: assignment switched to a spare, but the outgoing carrier stopped on item 3's only shelf during its return home. The right to commission the spare was valid; its remaining route could not obtain the item.

\paragraph{Fixed source and guard.}
Keep the 111 previous authored schedules, the physical simulator, shelf and notice laws, two recovery arms, and the 200-tick horizon. The actual stop times are unchanged. Both arms rerun native grid A* on a home-bound move with cells of shelves supplying any as-yet-unsettled class removed from the search, using the station's actual prefix at that cut. The vehicle can leave a shelf it already occupies; it cannot enter another such shelf while returning. If there is no permissible route, it waits. All unit movement and collision checks used earlier still apply. No second shelf, component, communicated fact, authorization, or future knowledge enters the control law. Interlocking against the exact station prefix presumes the local simulator-visible read available in both arms, not a proof that a deployed robot receives that read on time.

This rule rules out a specific source cut. Once an outgoing vehicle completes a protected move, its new position is not a shelf of any unresolved class at that cut. As long as it waits or continues protected movements, and the order has the same monotone settlement prefix, a stop can neither put it on a previously excluded shelf nor reopen the demand for an already settled class. This does not protect a vehicle that fails before its first legal move, fails during a required pick, changes from the protected motion mode, or faces a new governing order. Nor does the rule by itself authorize leaving an occupied resource, guarantee that such a detour exists, or certify that some other necessary route stays open.

\paragraph{Paired results.}
For the second physical stop at tick 23, the native discard plan with no guard cannot deliver when the first notice delay is 0 or 1. In each unguarded discard run at either delay the stop is on item 3's shelf, irrespective of the second report. With this rule the discard plan finishes at tick 93 or 94 respectively, including when the second report never arrives. The earlier amendment had moved the whole order duty to the spare; the lost robot's physical obstruction, not a missing grant, was what remained. At immediate notice, the second stop moves from cell $(4,4)$ to $(6,2)$; at one tick of delay the protected stop is at $(5,3)$. The repaired runs use nine new class pickups and 93 total unit moves; the unguarded runs spend six picks and 81 moves and never obtain item 3. These are more completed work and a different resource expenditure, not a claim of lower raw costs for a matched successful run.

The same guard on retaining loads reduces the immediate, timely-second-notice completion from 81 to 75, but still cannot settle when the second notice is completely erased. The affected carrier remains a commissioned supplier in that policy. More space to maneuver does not change its authority or fulfill its delivery duty. Under either routing law the early tick-20 fault does not finish in any of the tested schedules. In total ten unguarded noncompletions become native completions; no previously completing row becomes a noncompletion. The 111 rows and 812 actual home moves were independently reconstructed from simulator events. Those counts are deterministic development comparisons, not field frequencies.

This establishes a positive service envelope for the guarded discard rule at tick 100 in the declared two-delay family, with no second physical fault or with the tick-23 fault even when its notice is erased. Releasing a failed carrier's duty and making its withdrawal harmless to the spare are distinct conditions. A direct shelf-aware route planner can use the same simple guard, so this is not an advantage over the strongest established assembly or held-out I1/I2 acceptance. It provides an executable obligation to preserve a source during actual old-role withdrawal.
\end{document}
